\documentclass[11pt]{article}
\usepackage[T1]{fontenc}
\usepackage{mathpazo}
\usepackage[margin=1.15in]{geometry}
\usepackage{amsmath,amssymb,amsthm}
\usepackage{microtype}
\usepackage[hidelinks]{hyperref}
\usepackage{setspace}
\setstretch{1.12}

\newtheorem{definition}{Definition}
\newtheorem{proposition}{Proposition}
\newtheorem{corollary}{Corollary}
\theoremstyle{remark}
\newtheorem{remark}{Remark}

\newcommand{\Adm}{\mathrm{Adm}}
\newcommand{\E}{\mathbb{E}}
\newcommand{\Var}{\mathrm{Var}}

\title{Consistency Without Contribution:\\Fees, Selection, and the Accounting of Work}
\author{Flyxion\\Independent Researcher}
\date{October 2026}

\begin{document}
\maketitle

\begin{abstract}
\noindent A distributed ledger enforces consistency with a recorded history, and consistency is a different property from usefulness. This essay separates three tests that discussions of cryptocurrency routinely run together: a test of cash conservation, a test of productive benefit, and a test of responsibility after an operator exits. Productivity is defined as the net reduction in the resource cost of delivering a bundle of needs at a stated level of assurance, measured against the cheapest accessible alternative that does not use the system and counting costs over every party affected. Three conditional results are then proved. First, for a specified group of participants and a specified cash pool, the aggregate net cash position of the group, counting recoverable balances, equals minus the costs it pays to outsiders, so internal transfers cannot create a cash surplus; this is a statement about cash, and it says nothing by itself about trading profit, social waste, or who wins. Second, price information has no value to a decision problem when the state that matters is independent of the price, and it has value only if it changes the choice, so whether a market's prices are useless to everyone depends on a closure condition that has to be argued case by case. Third, selection among admissible states is neutral with respect to productivity when fitness and productivity are statistically independent under the reference measure, and is negative when the two depend on expenditure in opposite monotone directions. Published evidence on retail losses, wash trading, mining energy and electronic waste is used to ask where the premises hold. The argument is conditional: the energy objection depends on the consensus mechanism and not on the ledger, and the claims that survive that concession concern closure and the accounting of costs. Four cases from a recent documentary and from court reporting are examined, and a final section develops what survives the exit of an operator, including the observation that portfolio diversification lets investors tolerate failures whose costs remain concentrated on the people who depend on individual ventures.
\end{abstract}

\newpage

\section{Consistency and Usefulness}\label{sec:intro}

A ledger that cannot be silently altered is a real invention. The append-only chain of blocks described by Nakamoto \cite{nakamoto2008} allows parties with no prior relationship to maintain a shared history without a trusted intermediary, and the property it secures is consistency: every present claim that the protocol represents agrees with every recorded transfer that produced it. Consistency protects the recorded history and does not establish that the history is complete, since claims, obligations and contributions that were never represented in the protocol are outside what it can check. The question taken up here begins after that achievement is granted. Consistency with a log says nothing about whether the events in the log did anything for anyone, and a system can be flawlessly consistent while the activity it records only moves claims from one holder to another. The ledger certifies the movement. It does not certify that the movement was worth the resources used to perform it.

The gap can be put in a phrase: a ledger can be admissible while the value is absent. The phrase is easy to assert and hard to make exact, because the obvious ways of making it exact fail. A measure of work that is never given a concrete form makes the central proposition true by definition. A claim that redistribution cannot reduce work ignores the reply that settlement, custody and liquidity are services whose provision can lower the cost of something. A claim that speculative markets are lotteries conflates a zero-sum market with a negative-expectation game. A claim that the probability of productive configurations tends to zero needs an argument that the usual premises do not supply. And an argument from energy use has to allow for proof of stake, which is the first objection any informed reader raises. The essay therefore proceeds by definitions that can fail, each with its premise stated openly, and it keeps the structural idea that consistency with a log is a different property from productivity.

The thesis, in the form that survives scrutiny, has three parts. A system is productive only if it reduces the cost of delivering some bundle of needs that exist apart from the system, relative to the cheapest accessible alternative and counting every party who bears a cost. A market whose assets confer no claim on production, and whose prices improve no decision about production, allocation or risk-sharing in the population considered, is closed with respect to that population, and in a closed market the resources spent on information and on competition are lost without compensating gain. And the dynamics of such a market are driven by price and not by productivity, which makes them neutral toward productive configurations when the two are statistically independent and biased against them when fitness is purchased with expenditure that adds nothing. None of these parts is a statement about technology. Each is a statement about the structure within which a technology is used, and each can be tested. Alongside them runs a discipline that the later sections try to keep: a cash identity, a productive-benefit inequality, and the question of who answers for the consequences after an operator leaves are three separate tests, and passing one says very little about passing the others.

\section{Productivity as Net Work Reduction}\label{sec:productivity}

A usable definition has to be more concrete than a functional over an abstract field evolving in time. Fix a population $P$ and a set $N$ of needs, by which is meant capacities that members of $P$ must maintain whatever system they use: the ability to pay and be paid, to store claims safely, to insure against loss, to allocate capital to uses, to prove who owns what. The set $N$ is specified without reference to any particular system $\Sigma$, and this requirement does the real work in the definition, because without it a system could manufacture its own demand and then claim to have served it.

The counterfactual has to be stated with care, because it is where a favorable verdict is easiest to manufacture. Summing the cheapest way of meeting each need separately can overstate the cost of meeting them together, since an alternative institution may provide settlement, custody and verification through shared infrastructure, and a ledger compared with an artificially fragmented alternative will look productive for no good reason. The comparison is therefore made on a bundle of needs delivered jointly, at a stated level of quality, coverage, reliability, timing and volume, by arrangements that the population studied can actually reach.

\begin{definition}[Counterfactual cost of a bundle]\label{def:cost}
Let $S\subseteq N$ be a bundle of needs and let $q$ denote a level of assurance, meaning the quality, coverage, reliability, timing and volume at which the bundle is delivered. Let $\mathcal F_t^{-\Sigma}$ be the set of arrangements $\Gamma$ that are available at period $t$ to members of $P$ and that do not involve $\Sigma$, and let $R_t(\Gamma)$ be the total resource cost of operating $\Gamma$, expressed in a fixed numeraire such as labor-hours or energy-equivalents and counted over the domain $\Omega_{\mathrm{global}}$ of all parties materially affected, including costs borne by non-participants. Then
\[
C_t^{-\Sigma}(S;q)\;=\;\inf_{\Gamma\in\mathcal F_t^{-\Sigma}}\bigl\{\,R_t(\Gamma)\;:\;\Gamma\text{ delivers }S\text{ at assurance }q\,\bigr\}.
\]
When no accessible arrangement delivers the bundle at that assurance, the matched-service comparison is infeasible and is reported as such. Assigning a cost to going without would add an evaluative assumption that a purely resource-based comparison does not make.
\end{definition}

\begin{definition}[Net work reduction]\label{def:delta}
Let $s_t(\Sigma)\subseteq N$ be the bundle that $\Sigma$ actually delivers in period $t$, let $q_t$ be the assurance at which it delivers it, and let $R_t(\Sigma)$ be the total resource cost of operating $\Sigma$ over $\Omega_{\mathrm{global}}$. The net work reduction of $\Sigma$ in period $t$ is
\[
\Delta_t(\Sigma)\;=\;C_t^{-\Sigma}\bigl(s_t(\Sigma);q_t\bigr)\;-\;R_t(\Sigma).
\]
The system is \emph{positively productive over $[0,T]$} if $\sum_{t\le T}\Delta_t(\Sigma)>0$.
\end{definition}

The definition has an observable counterfactual built into it. The cost of delivering the same bundle without the system is the first term, the cost with it is the second, and productivity is the difference. Three features deserve comment. The domain $\Omega_{\mathrm{global}}$ matters, because a system that appears to save work among its participants may merely move cost to people who never joined, and the definition forbids that bookkeeping. The alternatives must be accessible to the people studied: a banking service that is cheap in principle and unavailable to the relevant users is not their counterfactual. And the numeraire conversion is a convention that will always be contestable for environmental and health burdens, which is a reason to report results under more than one conversion rather than a reason to drop the burdens. What the definition does not require is a welfare judgment. A person may enjoy a game of chance, and that enjoyment is real, but it does not lower the cost of delivering any bundle in $N$ and therefore does not register as productivity in this sense. The functional measures one thing, and the essay claims only what that thing can support.

The criterion is best described as comparative resource efficiency. A system that delivers an essential service at exactly the resource cost of its best accessible alternative receives $\Delta_t=0$, although it plainly contributes something useful, and that is acceptable provided two verdicts are kept apart: no demonstrated resource saving, which is what $\Delta_t\le 0$ reports, and no contribution, which it does not.

Two accounts have to be kept apart throughout. The first is financial: who paid what to whom, in money. The second is the account of resources: the labor, energy, hardware, administration and other inputs that were used up. A fee paid by a customer to an exchange is a cost in the first account for the customer and income for the recipient, so across $\Omega_{\mathrm{global}}$ it is a transfer, and it enters the second account only to the extent that it is evidence of the resources the recipient used in providing the service, under an estimation method that has to be stated. The cost $R_t$ in the definitions above is a cost in the second account, and the cash identity of Section~\ref{sec:cash} belongs to the first.

The definition has a corollary about appearances that deserves a concrete illustration, because a naive reading of productivity as visible output would get it wrong. A bag of potato chips is mostly air, and the observer who infers waste from the empty volume has measured the wrong quantity. The gas, usually nitrogen, cushions the contents against breakage in handling and slows oxidation, so that a smaller bag with less gas would raise breakage and spoilage and increase the work needed to deliver a given quantity of edible product. The comparison that matters concerns the delivery of intact, adequately preserved food at a stated assurance, and not the visible quantity of packaging. What looks unproductive when inspected locally contributes to the function of the larger arrangement, and the verdict can be reached only by counting costs and services across the whole process from production to consumption. The same discipline applies in both directions. It forbids calling an apparently idle component waste without examining what it does, and it equally forbids calling an apparently busy component productive without examining what the activity displaces. Much of what follows is an exercise in applying that discipline to a ledger and to the markets built on it.

The first consequence concerns pure redistribution, stated with its premise in the open.

\begin{proposition}[Pure redistribution]\label{prop:redistribution}
Suppose that in period $t$ the system $\Sigma$ delivers no need in $N$, so that $s_t(\Sigma)=\varnothing$. Then $\Delta_t(\Sigma)=-R_t(\Sigma)\le 0$, with strict inequality whenever the system consumes any resource.
\end{proposition}

\begin{proof}
The null arrangement of providing nothing is assumed to be accessible and costless, so that $C_t^{-\Sigma}(\varnothing;q_t)=0$. The assumption is needed because it is not true that every arrangement delivering nothing costs nothing, and the infimum is taken over accessible arrangements. Hence $\Delta_t=-R_t$. Moreover $R_t>0$ whenever any resource is consumed anywhere in $\Omega_{\mathrm{global}}$.
\end{proof}

The proposition is trivial, and that is the point. The premise is the only substantive claim: that the activity under examination delivers no need in $N$. Whether that is true of a given use of a ledger is an empirical question, and it is answered differently for a remittance corridor than for the trading of a token that confers no claim on anything. The remainder of the essay is largely about what makes the premise true and how much of observed crypto activity satisfies it.

\section{Cash Conservation, Fees, and the Lottery}\label{sec:cash}

The most common description of speculative markets as zero-sum is correct before costs and incomplete after them, and the difference is the whole of the first quantitative result. The result is an accounting identity about cash and has to be stated with its boundary. Consider a group $G$ of participants and a cash pool, held by an intermediary or a contract, which is empty at the start and whose only inflows are deposits by members of $G$. Let $D_T$ be the cumulative deposits of $G$ into the pool through the terminal date $T$, $W_T$ the cumulative withdrawals by $G$, $F^{\mathrm{in}}_T$ the costs paid out of the pool to parties outside $G$ (exchange and network fees, custody charges deducted from balances), $F^{\mathrm{out}}_T$ the costs that members of $G$ pay directly from outside the pool (hardware, electricity, separately paid fees), and $B_T$ the \emph{recoverable} terminal cash, meaning the cash the pool could actually pay out at $T$, as distinct from the claims that members hold against it. Define the aggregate net cash position of the group as
\[
N_T\;=\;W_T+B_T-D_T-F^{\mathrm{out}}_T .
\]

\begin{proposition}[Conservation of net cash position]\label{prop:conservation}
Suppose that the pool is initially empty, that its only inflows are deposits by $G$, and that its only outflows are withdrawals by $G$ and the payments $F^{\mathrm{in}}_T$. Then
\[
N_T\;=\;-\bigl(F^{\mathrm{in}}_T+F^{\mathrm{out}}_T\bigr)\;\le\;0 .
\]
\end{proposition}

\begin{proof}
Cash held in the pool at $T$ equals deposits minus withdrawals minus outside payments, so $B_T=D_T-W_T-F^{\mathrm{in}}_T$. Substituting into the definition of $N_T$ gives $W_T+(D_T-W_T-F^{\mathrm{in}}_T)-D_T-F^{\mathrm{out}}_T=-F^{\mathrm{in}}_T-F^{\mathrm{out}}_T$, and both cost terms are nonnegative.
\end{proof}

The premise that does the work is that no one outside $G$ deposits into the pool; if outsiders do, their deposits are a further inflow that the identity must include. Within that boundary the proposition says that transfers inside the pool cannot create a cash surplus for the group, and it says no more than that. The payments $F^{\mathrm{in}}_T$ and $F^{\mathrm{out}}_T$ are financial payments, entries in the first account of Section~\ref{sec:productivity} and not measures of resource cost. The left-hand side is the group's net cash position through the terminal date, with recoverable balances counted as the group's. It is not conventional realized trading profit, because a participant can sell a token at a profit and leave the proceeds in the pool, where they sit in $B_T$. A terminal balance is likewise not necessarily a loss, since it may be fully redeemable. The identity does not show that the winners are early entrants, which is a distributional claim requiring evidence about who bought and who sold and when. And it does not establish that the activity is socially unproductive, because outsiders who receive the payments $F^{\mathrm{in}}_T$ and $F^{\mathrm{out}}_T$ may be supplying useful services in return, which is the question that the productivity test of Section~\ref{sec:productivity} has to settle separately. Payments for goods and settlement of obligations require a separate account of token transfers and delivered services; their relation to this cash pool must be specified, since purchases made with tokens need not move cash through the pool at all. The proposition therefore does not say that crypto is a loss. It says that the purely financial part of it cannot generate a cash surplus for the group that trades, and that the size of the nonfinancial part decides everything else.

Two bodies of evidence bear on how the losses are distributed in practice. A Bank for International Settlements bulletin by Cornelli, Doerr, Frost and Gambacorta, using data on retail use of more than two hundred exchange applications between 2015 and the end of 2022, found that a majority of app users in nearly all the economies sampled probably lost money on their bitcoin holdings, with trading activity rising after the collapses of Terra/Luna and of FTX as larger and more sophisticated holders sold and smaller retail holders bought \cite{bis2023}. The estimate rests on an assumption that a user invests on the day the application is downloaded, and the authors present it as a proxy, but the direction is consistent with a cash identity in which gains for some are funded by deposits from others. The same bulletin observes that the turmoil of 2022 had little discernible effect on broader financial conditions, which its authors describe as underlining the largely self-referential nature of crypto as an asset class. That observation bears on the closure question taken up in the next section. On the supply side of volume, Cong, Li, Tang and Yang developed statistical tests for fabricated trades and found that wash trading is widespread on unregulated exchanges, averaging over seventy percent of reported volume on the unregulated venues in their sample of twenty-nine, while the regulated exchanges in the sample showed the patterns expected of genuine trading \cite{cong2022}. A market in which most reported volume on a large class of venues is fabricated is a market in which a large share of the apparent liquidity and apparent interest that attracts entrants is itself manufactured, and the cash that entrants deposit is what allows earlier holders to convert.

The familiar comparison with a lottery needs care in light of this. A lottery is a game with a published payout fraction below one and an operator who keeps the difference by design. A speculative token market has no operator who sets that fraction, and its expected net cash position for the group is negative because of fees and resource costs rather than because of a declared house edge. The two structures resemble each other in the distribution of outcomes, in which a few large winners are sustained by many small losers, and they differ in the mechanism that produces it. The resemblance is useful for intuition about who wins, and it should not be pressed to the claim that the market is a lottery, since in a market the winners are determined by timing and information rather than by chance, and that difference is what makes the next section necessary.

\section{Information, Price Discovery, and Closure}\label{sec:information}

The standard defense of trading in financial markets is that it discovers prices, and prices guide the allocation of real resources. The defense is sound where it applies, and Hirshleifer showed in 1971 where its limits lie \cite{hirshleifer1971}. In a pure exchange economy, information that arrives before others have it can be profitably traded on, and the private return to acquiring it can be large, while the social return is zero, because the information changes who holds what and does nothing to the quantity of goods available. Resources spent on acquiring such information are therefore a deadweight loss that the market nevertheless rewards. The result belongs to the exchange economy it is proved in, and where information bears on production or on risk-sharing the conclusion has to be reexamined. The classical analysis of Grossman and Stiglitz makes a related point in the other direction: if information is costly and prices fully reveal it, no one has an incentive to acquire it, so equilibrium prices cannot be fully informative \cite{grossman1980}, and the microstructure model of Kyle shows how informed traders' profits are paid by uninformed liquidity traders \cite{kyle1985}. The consistent lesson is that the social value of price information is not measured by the profit of those who trade on it. It is measured by how much better the decisions are that use the price.

That observation can be made precise with the standard value of information for a specified decision problem. Let a decision maker choose an action $a$ to maximize expected utility $u(a,\theta)$, where $\theta$ is a state that matters to that decision, and let $p$ be a market price. Under the usual existence and integrability assumptions, the value of observing $p$ is
\[
V(p)\;=\;\E\!\Bigl[\max_a \E[u(a,\theta)\mid p]\Bigr]\;-\;\max_a \E[u(a,\theta)] .
\]

\begin{proposition}[Independence implies zero value]\label{prop:coupling}
If $\theta$ is independent of $p$, then $V(p)=0$. Conversely, $V(p)>0$ only if no single action is optimal for almost every observed price. Dependence between $\theta$ and $p$ is necessary for positive value and is not sufficient for it.
\end{proposition}

\begin{proof}
If $\theta$ is independent of $p$, then $\E[u(a,\theta)\mid p]=\E[u(a,\theta)]$ for every action, so the inner maximum does not depend on $p$ and the two terms are equal. For the second statement, suppose one action $a^{*}$ is optimal for almost every $p$. Then $\max_a\E[u(a,\theta)\mid p]=\E[u(a^{*},\theta)\mid p]$ almost surely, so the first term equals $\E[u(a^{*},\theta)]\le\max_a\E[u(a,\theta)]$, and since the first term is always at least as large as the second, $V(p)=0$. Dependence without any change of optimal action therefore yields no value.
\end{proof}

Three relations have to be kept apart. Statistical coupling is dependence between a price and some state. Beneficial information is dependence that changes a choice for the better. Harmful exposure is a channel through which a price move affects people without improving any decision they make, as when a fall in token prices reduces the spending of holders through a wealth effect. A market can have the first without the second, and the third is a separate matter of cost and not of information. The conclusion that a market's prices are socially worthless is therefore much stronger than the observation that no outside actor consumes them, since information can benefit someone inside the market through insurance or allocation even when no outsider uses it. The condition needed is that the prices improve no relevant productive, allocative or risk-sharing decision anywhere in the population considered. Call a market \emph{closed} with respect to that population when the condition holds, and \emph{open} to the extent that it fails.

For a listed company the share price is informative about the firm's prospects, and boards, lenders and rivals take decisions on that basis, so the market is open and the discovery of prices has a downstream consumer. For a token that confers no claim on cash flows, no right to a service, and no obligation of any counterparty, the price is informative about few states that matter to anyone's productive or allocative choices, the principal exceptions being the token's own use as a means of payment or collateral. For such a market the closure condition is plausible and is an empirical claim, not a theorem, and where it holds Proposition~\ref{prop:coupling} together with Hirshleifer's result gives a clear conclusion: the information the market produces improves no decision about production or allocation, the resources used to acquire it and to compete on it are not matched by a social return of that kind, and the rewards that attract those resources are transfers. The evidence cited above, that stress in the token market had a limited effect on broader financial conditions, is what a market with limited channels to the rest of the economy would show, and it is evidence about channels of propagation and not about whether the market is useful to those inside it.

The qualification matters and should not be skipped. Informational openness is not zero in general. Where tokens are used as a means of payment or accepted as collateral, their prices can improve decisions about whether to accept them, how much to lend against them and how to price contracts denominated in them, and to that extent the market is informationally open in proportion to the size of those uses. Financial interconnectedness is a different relation. Lending against tokens, bank deposits held for crypto firms, and wealth effects on holders who spend in the ordinary economy are channels through which shocks propagate, and they establish exposure and not beneficial information. A market can be informationally closed and financially connected, and neither relation follows automatically from the other. The claim of this section is that informational closure is a matter of degree, that for assets without claims on production the degree is high, and that the social value of price discovery should be assessed against the openness actually present and not assumed from the analogy with equity markets.

\section{Degeneracy as Selection Neutrality}\label{sec:selection}

Call a system degenerate when the set $\Adm(L)$ of states consistent with its log is large and the dynamics fail to select the productive subset $\Adm^{+}(L)$. It is tempting to add that the probability assigned to productive states then tends to zero, but that does not follow from the premises, and what does follow is more interesting. Suppose the state space is finite, and let $\mu$ be a probability measure on $\Adm(L)$ describing how many ways each configuration can arise. Let $f$ be the fitness the dynamics reward, for a market the price-appreciation proxy that order flow pushes upward, and let $g\in\{0,1\}$ indicate membership of $\Adm^{+}(L)$. Suppose the stationary distribution of the dynamics has the Gibbs form
\[
\pi_\beta(x)\;=\;\frac{e^{\beta f(x)}\,\mu(x)}{Z_\beta},\qquad Z_\beta=\sum_x e^{\beta f(x)}\mu(x),
\]
with $\beta\ge 0$ measuring how strongly fitness is selected. This is the stationary distribution of dynamics whose transition rates satisfy detailed balance with respect to it, as Metropolis-type rules with target $\pi_\beta$ do, and it should be read as a null model for markets driven by price. That nothing in the update rule mentions productivity is not enough for neutrality. An update rule that never refers to productivity can still select for it, or against it, through traits that are correlated with it, and the sufficient condition proved here is statistical independence under the reference measure, which is generally not necessary.

\begin{proposition}[Selection neutrality]\label{prop:neutral}
If $f$ and $g$ are independent under the probability measure $\mu$, then $\pi_\beta(\Adm^{+}(L))=\mu(\Adm^{+}(L))$ for every $\beta\ge 0$.
\end{proposition}

\begin{proof}
Independence under $\mu$ gives $\E_\mu[e^{\beta f}g]=\E_\mu[e^{\beta f}]\,\E_\mu[g]$. Dividing by $Z_\beta=\E_\mu[e^{\beta f}]$ yields $\E_{\pi_\beta}[g]=\E_\mu[g]$, which is the claim because $g$ is the indicator of $\Adm^{+}(L)$.
\end{proof}

Strong selection on price, then, does not by itself drive productive configurations out. If fitness and productivity are independent it leaves their weight at whatever the reference measure assigns, and a system can be degenerate in the weaker sense that it has no preference for what matters. The stronger conclusion requires an additional structure, and the mining race supplies it. In a competition with free entry for a prize of fixed value, entrants spend until the marginal return to spending is exhausted, so that aggregate expenditure tends toward the value of the prize. This is the classical logic of rent dissipation \cite{tullock1967,posner1975}, and the economics of proof of work has been analyzed in this way by Budish and by Auer, who relate the cost of securing the chain to the rewards it pays \cite{budish2018,auer2019}. Let $r:\Adm(L)\to\mathbb R$ denote the resource expenditure embodied in a state.

\begin{corollary}[Negative selection under expenditure]\label{cor:negative}
Suppose $f=F\circ r$ and $g=G\circ r$, where $F$ is nondecreasing and $G$ is nonincreasing. Then $\E_{\pi_\beta}[g]\le \E_\mu[g]$ for every $\beta\ge 0$.
\end{corollary}

\begin{proof}
The function $x\mapsto e^{\beta F(r(x))}$ is nondecreasing in $r$ and $x\mapsto G(r(x))$ is nonincreasing in $r$, so by Chebyshev's association inequality their covariance under $\mu$ is at most zero. Hence $\E_\mu[e^{\beta f}g]\le \E_\mu[e^{\beta f}]\,\E_\mu[g]$, and dividing by $Z_\beta=\E_\mu[e^{\beta f}]$ gives the claim.
\end{proof}

The functional form $f=F\circ r$ and $g=G\circ r$ is stated explicitly because variation among states with equal expenditure could otherwise defeat the covariance argument. The premises are also exactly what a defender of the system would contest, since the second one says that additional expenditure never makes a state more productive. That premise can be examined case by case, and it is where the distinction between consensus mechanisms and between kinds of computation enters. If extra expenditure completes tasks that are independently required, raises assurance that cannot be had more cheaply, or displaces expenditure on heating, then productivity does not fall with expenditure and the corollary does not apply. The proposition and its corollary stand in place of a claim of vanishing probability with two claims that can be falsified: neutrality when fitness and productivity are independent, and negative selection when both depend on expenditure in the stated monotone directions.

\section{Consensus Mechanisms and the Material Ledger}\label{sec:material}

The material costs of the system depend on how agreement on the ledger is reached, and a statement that cryptocurrency is energy-intensive is a statement about proof of work and not about ledgers. On the Bitcoin network, estimates of annual electricity consumption have been broad. When the Cambridge Centre for Alternative Finance launched its electricity consumption index in 2019, its central estimate was in the neighborhood of fifty terawatt-hours a year within a band running from about twenty-one to about a hundred and forty-six, and its authors stressed that reliable estimates are rare and diverge widely between models \cite{cbeci2019}. Current figures are published by the index and should be taken from it directly. What is structurally important is the dependence of expenditure on price. Because the reward for adding a block is denominated in the token, free entry drives spending toward the value of that reward, so electricity use tracks the price of the token and not the number or the usefulness of the payments it settles.

Proof-of-work expenditure is nonetheless not nothing in exchange, and the accounting has to say what it buys. It supplies a security service, since the cost of attacking the chain rises with the resources committed to it, and the economics literature cited above ties that security to the rewards paid \cite{budish2018,auer2019}. The correct criticism therefore compares the resources required for equivalent assurance, and it does not classify all puzzle computation as serving no need. To the extent that the expenditure buys assurance that cannot be had more cheaply, it is the cost of a service. To the extent that equivalent assurance can be delivered for less, the difference is waste relative to the alternative, and it enters $\Delta_t$ through the counterfactual cost of Definition~\ref{def:cost}. Whether Corollary~\ref{cor:negative} applies depends on whether additional expenditure buys additional assurance or only offsets the expenditure of competitors, and a rise in price that raises expenditure without any change in the service provided is the case in which it does.

The point has to be set against the transition of Ethereum from proof of work to proof of stake in September 2022, which according to the Ethereum Foundation reduced the energy consumption of the network by about $99.95$ percent \cite{ethereum2022}. The figure is self-reported by the developers and has not been verified independently here, though the direction and rough magnitude are not disputed in the sources consulted. It follows that the energy cost is a property of a consensus design and not of the ledger or of tokens as such, and that any argument for the structural unproductivity of token markets cannot rest on electricity. A case for intrinsic failure that rested part of its weight on energy would be wrong on this point. What remains after the concession is the closure argument of Section~\ref{sec:information} and the neutrality result of Section~\ref{sec:selection}, neither of which depends on the consensus mechanism, together with the cash identity of Section~\ref{sec:cash}, which is independent of it as well. Proof of stake changes $R_t(\Sigma)$, and through it the sign of $\Delta_t$ in the settlement use case; it does not make a token with no claim on production into an open market.

Hardware is a second material channel and a clearer one. De Vries and Stoll estimated that Bitcoin mining generated about thirty thousand seven hundred metric tons of electronic waste annually as of May 2021, comparable to the small information-technology equipment waste of the Netherlands, using an assumed average device lifetime of about a year and a quarter, and they warned that demand for mining hardware could disturb the semiconductor supply chain \cite{devries2021}. The figure is a model-based estimate whose level depends on the price assumptions of its date, and the per-transaction figure sometimes quoted from the same study, in the neighborhood of two hundred and seventy grams, divides an expenditure that scales with price by a throughput that does not, which makes it a convention of allocation rather than a physical coefficient. Both caveats reduce the precision of the claim without removing its sign. The waste is real, it is generated where the hardware is disposed of, and the people who bear the associated exposure are not the participants of the market. This is the sense in which the domain $\Omega_{\mathrm{global}}$ in Definition~\ref{def:cost} is not decorative. Where hardware races are part of the mechanism, the externalities enter $R_t$ and should be reported, with the uncertainty of the estimates, rather than ignored or exaggerated.

\section{Where the Inequality Could Hold}\label{sec:where}

An argument that proved unproductivity everywhere would prove too much, and the framework is built so that it does not. The inequality $\Delta_t(\Sigma)>0$ is a claim about cost, and for each use of a ledger it can be checked. Consider a cross-border payment corridor in which the existing alternatives are slow, costly and exclusionary. The bundle is the ability to move value across a border at stated speed and reliability, the counterfactual cost $C_t^{-\Sigma}(S;q)$ is the cost of the best accessible alternative including its fees and delays, and the cost of the ledger is its resource cost, namely its settlement-layer expenditure including any external burdens and the on- and off-ramps through which users enter and leave, with fees entering as evidence of that cost only under a stated estimation method. If a proof-of-stake settlement layer delivers the bundle at a cost below that of the accessible alternative, the inequality holds, and nothing in this essay says otherwise. Consider a registry of credentials or of provenance, where the need is the ability to verify a claim cheaply by a party with no relationship to the issuer. The counterfactual is a trusted intermediary, and the ledger is productive if its cost is below the cost of maintaining and auditing that intermediary at equivalent assurance. In both cases the events recorded correspond to transformations of the world that exist whether or not the ledger records them, the information in the ledger is consumed by decisions outside it, and the market is open in the sense of Section~\ref{sec:information} by construction.

By contrast, a token that confers no claim on production, trades on markets whose volume is largely fabricated on the unregulated venues studied, funds the cash positions of its early sellers from the deposits of later entrants, and consumes resources in a competition for a reward denominated in itself, falls in the case that Proposition~\ref{prop:redistribution} describes, to the extent that it delivers no need in $N$. It is useful to note that the classification is not made by the technology. The same ledger supports both uses, and it follows that the relevant unit of analysis is the use and its openness, not the technology and its reputation. The fraction of activity that is of the first kind is an empirical matter on which this essay offers no figure, because the sources assembled here do not supply one, and an honest version of the claim is that the burden of measurement falls on those who assert that the productive fraction is large.

\section{Four Cases from the Record}\label{sec:cases}

The abstract results can be set against a recent documentary account of the same market. In \emph{Everyone Is Lying to You for Money}, released by CBC Docs in 2026, the actor and author Ben McKenzie follows crypto promotion from celebrity advertising to its collapses \cite{mckenzie2026}. It is journalism and advocacy rather than data, and its figures are used below only as the film's own reporting or as interview statements, each marked as such, with court outcomes taken from news reports of the proceedings. The material is illustrative and not probative: it shows how the premises of the preceding sections look in particular cases, and it does not estimate how common those cases are. The cases are useful because each exposes a different premise of the preceding sections.

The first case concerns the scale at which a claimed need was delivered. El Salvador made Bitcoin legal tender in 2021 with remittances as the headline use, the need being cheaper transfer of money from workers abroad to families at home, a need that exists independently of any token and for which the counterfactual cost is that of the accessible channels the families already use. The film reports that, by the government's own figures, only about $1.6$ percent of remittances were paid in Bitcoin, and in an interview with McKenzie, who put it to him that the remittance use was not working, Sam Bankman-Fried agreed that it was not working at that time and said that the majority of people were using crypto as a financial asset rather than as a payment vehicle. A low share of transactions is not an empty bundle. In the vocabulary of Definition~\ref{def:delta}, $s_t(\Sigma)$ records which needs are delivered and not how many people receive the service, and a system that handles a small share of remittances can still deliver the same kind of need to its users. Transaction volume and population coverage belong in the specification of the service, and low uptake limits the scale of demonstrated provision, so whether total savings were positive still requires costs and volumes that the film does not give. The film also reports that the planned airport for the proposed Bitcoin City would displace a fishing community offered a fraction of the value of its land, which is an externality borne by non-participants and so a term of $R_t(\Sigma)$ over $\Omega_{\mathrm{global}}$ in exactly the sense required.

The second case is a closed loop in which the cash identity and the closure condition can be seen operating. Celsius marketed itself as a replacement for banks, and its former chief executive Alex Mashinsky pleaded guilty in December 2024 to commodities fraud and to a scheme to manipulate the price of the firm's token, and was sentenced in May 2025 to twelve years in prison \cite{nbc2025,decrypt2025}. Reporting on the case describes the charges as misappropriating customers' funds and manipulating the price of the token, and the company's customers were reported to have lost more than five billion dollars \cite{decrypt2025}. The loss in such a case is not the terminal cash balance of Proposition~\ref{prop:conservation}, which may be fully recoverable in other settings. It is the shortfall between the claims that customers hold against the firm and the assets that can actually be recovered, and the distinction between outstanding claims and recoverable assets is the one that identity forces. The film records Mashinsky, in an audio-only exchange after the video had stopped, estimating that perhaps ten or fifteen percent of the money in the crypto system was real money that had gone in and that everything else was bubble, and McKenzie repeats the figure to Bankman-Fried. That is an offhand estimate by an interested party and carries no evidential weight beyond its source, but its shape matches the cash identity, according to which the cash that can be withdrawn is bounded by the cash that went in.

The third case, FTX, concerns financial interconnectedness and the propagation of failure, which is a different relation from the informational openness of Section~\ref{sec:information}. Bankman-Fried was found guilty on all counts after a month-long trial and sentenced to twenty-five years in prison in March 2024 \cite{nbc2025,mckenzie2026}. The film's account is that two lines of the exchange's code prevented users from overdrawing their accounts and that he had an employee add a line allowing the firm to borrow as much as it wanted from customers, and its larger argument concerns spread: that integrating crypto with the regulated banking system risks transmitting its failures to depositors who never chose the risk. The film cites the failure of three United States banks within five days in March 2023, which it describes as critical counterparties of the crypto industry, and notes that the government guaranteed customers' deposits. The causal story is the film's and is not checked independently here, and it should be kept apart from the earlier evidence, since the BIS bulletin discussed above, which predates those failures, found little discernible spillover from the 2022 turmoil. The two statements are compatible: a market can be informationally closed and financially connected through particular channels, and the bank counterparties are such a channel. A market with few channels can still impose heavy costs on its own participants and, through other channels such as the externalities of Section~\ref{sec:material}, on people who never took part in it, so the absence of spillover into financial conditions is not evidence that the market was harmless.

The fourth case concerns the quality of the volume figures that attract participants. An interviewee in the film, a former video-game cheater who says he has been exposing crypto fraud since 2017, describes how one person trading between two accounts can walk a price upward and asserts that between ninety and ninety-five percent of all crypto trading is of that kind. The peer-reviewed estimate discussed in Section~\ref{sec:cash} is more modest and better supported, with wash trading averaging over seventy percent of reported volume on the unregulated exchanges in a sample of twenty-nine and absent from regulated ones \cite{cong2022}. The more modest figure should be preferred, since it comes with a method and a sample, and it should be reported as such in any argument that rests on it. The difference illustrates the discipline of the essay in a small way: a claim does not become more reliable by being more dramatic, and the structural argument does not need the larger number.

\section{What Survives an Exit}\label{sec:exit}

The cases share a feature that the accounting of Section~\ref{sec:productivity} does not yet capture. When a firm collapses or its principals leave with the proceeds, what remains is a set of obligations, a lost record of who contributed what, and a population left holding claims whose value depended on the continued operation of the thing that has just ended. The ledger records faithfully what was entered in it, and consistency protects that history without establishing that the history is complete. Obligations and contributions that were never represented are exactly what it cannot show, and even a faithful record of what was entered is little comfort to those who are owed. A constructive alternative has to say what survives an exit and who remains answerable, and the three tests named at the outset, cash conservation, productive benefit and responsibility after exit, are the way to organize the question.

One line of proposals, developed in a separate essay on contribution, identity and governance after tokenization \cite{flyxion2026ledger}, begins from the observation that a single token typically bundles three functions that need not go together: a record of who did what, a medium of payment, and a source of authority over decisions. When the three are fused, the price of the token becomes the fitness function $f$ of Section~\ref{sec:selection}, and payment and authority are pulled along with it, so that early participation converts into perpetual control and a profitable exit can erase the record of who did the work. Unbundling treats each function as a separate question requiring its own justification. The record keeps contributions with their actors, circumstances, evidence, attestations, objections and scope, and does not reduce them to interchangeable units of price. Payment remains possible by ordinary currency, mutual credit, grants, gifts or exchange in kind, and records a particular agreement rather than a universal measure of the contributor's worth. Standing in a decision depends on a person's relation to that decision, whether affected by it, responsible for carrying it out, or competent to assess the evidence, and these relations are kept distinct, so that expertise supplies evidential weight without sovereignty and being affected supplies a claim to consideration without token ownership. Historical contribution stays on record, but its bearing on present authority has to be argued again. Technically, such a record needs only signed append-only logs, content-addressed evidence and replication, which a version-control object model already provides, and does not need a public blockchain with a native token. The description here is of the structure of the proposal and not of its details, which belong to the separate essay.

Two features of this proposal connect to the argument of the present essay. First, separating the record from the price removes the coupling between fitness and balance, so that the evidence of work can be made the object of scrutiny and selection pressure can be applied to the productive subset directly, where the neutrality of Proposition~\ref{prop:neutral} would otherwise leave it unexamined. Second, the treatment of useful computation in that proposal is an application of Definition~\ref{def:cost}, with one correction to the way the question is usually posed. A protocol can verify that a computation was carried out correctly, and that is a separate matter from whether the result serves a purpose. Puzzle computation is not without function, since it supplies the security service discussed in Section~\ref{sec:material}, and the question is therefore whether useful computation can meet the relevant security requirements while also providing additional, independently justified benefits, after verification, duplication and other overhead are counted. Computation directed at a task that an independently existing scientific, ecological or infrastructural need already requires counts only if the need is specified without reference to the reward mechanism, since otherwise useful work becomes a name for whatever the mechanism induces. Placing computation where its heat is wanted lowers $R_t$ only to the extent that the recovered heat displaces another source, so the credit has to be computed against the actual alternative and not assumed. Finally, useful work does not have to mint a transferable asset or confer authority. It can justify a contractual payment and an acknowledgment, with authority left to be authorized separately.

The record and the unbundling do not by themselves supply a mechanism for responsible closure, and the structure of investor-funded ventures shows why one is needed. A diversified investor holds claims on many ventures, each of which may fail with substantial probability, and the portfolio is rational for the investor because the successes can repay the failures. The tolerance that results belongs to the investor's position and not to the ventures. A simple model makes the asymmetry exact. Suppose that $J$ ventures fail independently with the same probability $q$, and let $X_j\in\{0,1\}$ indicate failure of venture $j$. Suppose further that each venture receives an equal stake, that a failed venture returns nothing, and that a successful one returns a multiple $m>1$ of its stake, so that the multiple earned by venture $j$ is $M_j=m(1-X_j)$. The investor's net return is $\Pi_J=\tfrac1J\sum_j M_j-1$. The average $\bar X=\tfrac1J\sum_j X_j$ is the gross principal-loss fraction attributable to failures: it assumes equal stakes, complete loss on failure and no loss of principal on success, and it leaves out the gains of the successful ventures, which are what $\Pi_J$ adds back.

\begin{proposition}[Diversification and dependence]\label{prop:portfolio}
Under these assumptions, $\E[\bar X]=q$ and $\Var(\bar X)=q(1-q)/J$. The investor's net return satisfies $\E[\Pi_J]=(1-q)m-1$, which is positive exactly when $(1-q)m>1$, and $\Var(\Pi_J)=m^2q(1-q)/J$. A person whose livelihood depends on one particular venture $j$ suffers the loss $X_j$, with $\E[X_j]=q$ and $\Var(X_j)=q(1-q)$. If the failure indicators have a common pairwise correlation $\rho$, then $\Var(\bar X)=q(1-q)\,\bigl(1+(J-1)\rho\bigr)/J$ and $\Var(\Pi_J)=m^2\Var(\bar X)$.
\end{proposition}

\begin{proof}
The means follow from $\E[X_j]=q$. Since $\Pi_J=m(1-\bar X)-1$, we have $\E[\Pi_J]=m(1-q)-1$ and $\Var(\Pi_J)=m^2\Var(\bar X)$. For independent indicators the variance of an average of $J$ terms is $q(1-q)/J$. With common correlation $\rho$ the variance of the sum is $Jq(1-q)+J(J-1)\rho\,q(1-q)$, and dividing by $J^2$ gives the stated expression.
\end{proof}

The model expresses how frequent failure can coexist with investor profitability: when $(1-q)m>1$ a portfolio in which most ventures fail still earns a positive expected return, and a large multiple on rare successes lets one exceptional success repay many failures. Diversification then reduces the dispersion of the investor's loss and of the investor's return and leaves the failure frequency where it was, and when failures are correlated, as they will be to the extent that ventures share a price cycle, the reduction is smaller and does not vanish as $J$ grows unless $\rho=0$. The dependent faces the full variance of a single venture, because employment, supply contracts, local services and the claims of customers are tied to the particular venture and cannot be spread across the portfolio. The investor's return measures the portfolio, and it neither measures the losses of the dependents nor compensates them. The model also supports a caution in the other direction. Failed experimentation can generate knowledge and reusable capacity, so evaluating every failed venture as pure waste would be a mistake, and the appropriate comparison covers the whole program at once. Summing counterfactual costs venture by venture would reintroduce the aggregation problem of Section~\ref{sec:productivity}, since the outputs of several ventures may be cheaper to provide jointly and what is learned from one may overlap what is learned from another. The program is therefore treated as a single system $\Sigma_{\mathrm{prog}}$ delivering a single bundle $S_{\mathrm{prog}}$, which includes the retained knowledge and reusable capacity at a specified assurance, so that nothing is credited twice:
\[
\Delta^{\mathrm{prog}}_{[0,T]}\;=\;C^{-\Sigma}_{[0,T]}\bigl(S_{\mathrm{prog}};q\bigr)\;-\;R_{[0,T]}\bigl(\Sigma_{\mathrm{prog}}\bigr)\;-\;K_T ,
\]
where $R_{[0,T]}(\Sigma_{\mathrm{prog}})$ is the total resource cost of the program over $\Omega_{\mathrm{global}}$, and $K_T$ is the additional resource cost of completing responsible closure, excluding resources already counted in $R$. A single venture is the special case of one component. If retained learning is instead entered as a separate term, it must measure incremental benefits that appear in no other term, including benefits already attributed to other ventures. In either form the benefit of retained learning is zero if nothing transfers, and like every other term it needs evidence and cannot be assumed.

The term $K_T$ is a resource cost, and it belongs to the second account of Section~\ref{sec:productivity}. Repayment of an outstanding monetary claim is a transfer and belongs to the first; the labor, administration, relocation, restoration and other resources needed to wind a venture down responsibly are resource costs, and both matter, though they enter different accounts. Recognizing a cost is also not the same as financing it. Whether resources are actually reserved to meet residual obligations is a separate institutional requirement, in the same way that preserving a contribution on a record does not pay its contributor. In practice $K_T$ has at least three components. Physical and environmental remediation covers decommissioning, site restoration and the handling of electronic waste for systems with a physical footprint. Archival cost covers preserving the historical state and records that users and dependents need in order to verify past claims after the operators have gone, which a ledger that nobody maintains does not supply by itself. Transition cost covers the resources needed to move dependents to an alternative arrangement, such as relocation, retraining, and the migration of users' records and obligations. A liquidity buffer that backs monetary claims is not part of $K_T$, because it is a transfer and belongs to the first account, although it may have to be held alongside the reserve. Limited liability sets a floor under the investor's loss at the amount staked and leaves the remaining costs to whoever depends on the venture, which is why the question of who funds $K_T$ is answered neither by the ledger nor by the register of owners.

One simple mechanism is a mandatory reserve. Suppose that each venture requires, besides a unit stake, a reserve $K>0$ placed in escrow outside the control of any token holders, so that governance cannot vote it away, which is the fusion of record, payment and authority that the unbundling above is meant to prevent. Suppose the reserve is returned to the investor if the venture succeeds and is applied in full to closure, so that $K$ is the monetary amount sufficient to procure the closure resources $K_T$ at specified prices, if it fails; the reserve finances closure and does not reduce its resource cost. Let a successful venture pay a gross amount $m>1$ and a failed one nothing, as in the model above, and let $R_j$ denote the investor's net monetary payoff on venture $j$ per unit of operating stake, which is not a rate of return on total committed capital, so that without a reserve $R_j=m(1-X_j)-1$ and the unconstrained portfolio return is $\Pi_J=\tfrac1J\sum_j R_j$.

\begin{proposition}[Mandatory closure reserve]\label{prop:reserve}
With the reserve, the outlay on each venture is $1+K$, and the net return is $R^{\mathrm{res}}_j=m-1$ on success and $R^{\mathrm{res}}_j=-(1+K)$ on failure. For independent failures with probability $q$, the portfolio return $\Pi^{\mathrm{res}}_J=\tfrac1J\sum_j R^{\mathrm{res}}_j$ satisfies
\[
\E\bigl[\Pi^{\mathrm{res}}_J\bigr]\;=\;(1-q)m-1-qK\;=\;\E[\Pi_J]-qK,
\qquad
\Var\bigl(\Pi^{\mathrm{res}}_J\bigr)\;=\;(m+K)^2\,\frac{q(1-q)}{J}.
\]
The reserved venture has positive expected net payoff, under these assumptions, exactly when $(1-q)m-qK>1$, and its expected return on total committed capital is $\bigl((1-q)m-1-qK\bigr)/(1+K)$.
\end{proposition}

\begin{proof}
On success the payout is $m+K$ against an outlay of $1+K$, giving $m-1$. On failure the payout is zero against the same outlay, giving $-(1+K)$. Hence $\E[R^{\mathrm{res}}_j]=(1-q)(m-1)-q(1+K)=(1-q)m-1-qK$, and linearity gives the mean of the average. The two outcomes differ by $(m-1)+(1+K)=m+K$, so each $R^{\mathrm{res}}_j$ has variance $(m+K)^2q(1-q)$, and the variance of an average of $J$ independent terms is that quantity divided by $J$.
\end{proof}

The reserve lowers the investor's expected net payoff by $qK$ per venture, which is the expected monetary cost of winding down failed ventures responsibly, so the mechanism internalizes that cost and does not create it. Ventures that were marginally profitable with the closure cost left to others fall below the threshold once it is pre-funded, which is the intended effect and also shows why the cost is otherwise externalized. Diversification still lowers the variance of the investor's return as $J$ grows, but it does not remove the expected outlay $qK$, so the cost of termination is borne by the capital structure and no longer by workers, suppliers, users and communities. The result is conditional on strong assumptions. It treats $K$ as known in advance and exactly sufficient, ignores the time value of money held in escrow, assumes that failure can be verified so that the reserve is released correctly, and holds the failure probability fixed although a reserve could raise it by making failure cheaper to the operator or lower it by making failure costly. It covers resource costs only and leaves the backing of monetary claims to a separate buffer. Each of these is an empirical or institutional question that would have to be answered before the mechanism could be relied on, and the proposition is offered as a first measure of what pre-funding costs and not as a design. What it adds is the statement of the problem in a form that can be audited.

The cash identity of Section~\ref{sec:cash} is an accounting constraint on a specified group taken in aggregate and is not a threshold of profitability for any individual investor, so a venture whose investors profit while its dependents absorb $K_T$ has produced a private financial gain and has left its productive balance and its closure responsibilities unresolved. Private returns, global resource savings and funded responsibility after exit are independently assessable, and the portfolio model gives a precise reason why success on the first cannot stand in for success on the other two.

\section{Objections}\label{sec:objections}

The first objection is that markets provide liquidity and that liquidity is a service. The reply is that liquidity delivers a need only for someone who needs to convert holdings into something else, and the question is again what the holdings are claims on. Liquidity in a market for claims on production lowers the cost of capital allocation. Liquidity in a market whose holdings are claims only on future buyers lowers the cost of entering and exiting a position, which is a service rendered to the speculators themselves and delivers no need in $N$. This is not a ruling that such liquidity is worthless to those who use it, since the functional is not a welfare function, and the claim here is limited to its effect on the cost of delivering bundles that the population needs independently.

The second objection is that a speculative asset can serve as a hedge or store of value, which are legitimate needs. The reply is conditional on the asset's payoff structure: an instrument provides risk-sharing only if its returns covary with states in which holders are made worse off in the opposite direction, and an asset whose price is driven by flows of attention and capital in the market itself provides it only incidentally. This is an empirical claim about covariances, and it can be checked, and where it is established the corresponding need is delivered and enters $s_t(\Sigma)$. The same applies to the possibility that a token serves as a hedge against failures of the monetary system in particular jurisdictions, a case in which the cost of the accessible alternatives is high and the productive claim is strongest.

The third objection is that people enjoy speculation and that enjoyment is a good. This is correct and is deliberately set aside by the construction of the functional. A consumption good that gives pleasure is not thereby a reduction in the cost of delivering a needed bundle, and the classification of speculative trading as a form of entertainment is compatible with every result above. What follows from the classification is that the language of technological progress and of economic empowerment, which claims a productive function, should be assessed under the productive criterion, while the entertainment claim should be assessed under the criteria that apply to other forms of gambling, including whether participants understand the odds and whether the odds are disclosed.

The fourth objection is that the test, applied evenhandedly, condemns a large part of conventional finance. It condemns exactly the part that fails the closure and cost tests, and that is a feature rather than a defect. Trading on short-horizon information about assets whose prices feed no allocation decision has been criticized on the same grounds in the conventional literature that begins with Hirshleifer, and the framework offers a way to ask how much of conventional trading is of that kind. The difference between the two domains is one of degree in openness and in regulation, and the data cited above on wash trading show a difference in the second respect: the regulated exchanges in the sample displayed the signatures of genuine trading and the unregulated ones did not.

The fifth objection is that proof of stake and other redesigns dissolve the problem. They dissolve the energy part of it and shrink the externality part, and they alter $R_t$ in the favorable direction. They do not alter openness, and they do not alter the cash identity of Proposition~\ref{prop:conservation}, in which the costs may fall but remain nonnegative, so that the group's net cash position remains nonpositive. A cheaper ledger lowers the cost of a closed market. It does not open it.

\section{Conclusion}\label{sec:conclusion}

The consistency of a ledger is a property of its record, and the usefulness of what it records is a property of the world. The essay has tried to make the gap between them measurable, and to keep three tests apart that are easily confused. The first is a cash test: for a specified group and pool, the group's net cash position equals minus the costs it pays to outsiders, which says that internal transfers create no surplus and says nothing yet about waste or about who gains. The second is a test of productive benefit: a system is productive when it delivers a bundle of independently existing needs at an assurance and a cost below those of the cheapest accessible alternative, counting every party affected, and the verdict on a market turns on whether its prices improve any decision about production, allocation or risk-sharing and on whether selection in it is neutral or negative with respect to productivity. The third is a test of responsibility after exit, which asks who bears the cost of failure when diversified investors tolerate it and dependents cannot. The empirical record assembled from published studies is consistent with the first two results for the speculative part of the market, in showing majority retail losses on the sample studied, in showing fabricated volume on unregulated venues, and in showing that stress in the token market did not propagate widely into broader financial conditions in 2022, a fact about channels and not about harm. It is also consistent with the exceptions the framework leaves open, since the energy objection fails once the consensus mechanism changes, and uses of ledgers in which the events recorded are consumed by decisions outside the ledger can satisfy the productivity inequality.

The remaining work is empirical and unglamorous. For each use of a ledger, the bundle of needs must be named, the accessible alternative must be priced at equivalent assurance, the cost of the ledger must be counted over its full domain with the uncertainty reported, the closure of the market must be argued and not assumed, and the cost of orderly closure must be shown to be met by someone. The framework will sometimes return a positive number. A theory that cannot return one, or a theory that returns one for every application whatever the facts, is not doing the work that a theory of value is for.

\begin{thebibliography}{99}

\bibitem{auer2019} R.~Auer, \emph{Beyond the doomsday economics of ``proof-of-work'' in cryptocurrencies}, BIS Working Paper No.~765, Bank for International Settlements, 2019.

\bibitem{bis2023} G.~Cornelli, S.~Doerr, J.~Frost, and L.~Gambacorta, \emph{Crypto shocks and retail losses}, BIS Bulletin No.~69, Bank for International Settlements, 20 February 2023.

\bibitem{budish2018} E.~Budish, \emph{The economic limits of Bitcoin and the blockchain}, NBER Working Paper No.~24717, 2018.

\bibitem{cbeci2019} Cambridge Centre for Alternative Finance, \emph{Cambridge Bitcoin Electricity Consumption Index}, University of Cambridge, launched 2019.

\bibitem{cong2022} L.~W. Cong, X.~Li, K.~Tang, and Y.~Yang, \emph{Crypto wash trading}, NBER Working Paper No.~30783, 2022.

\bibitem{decrypt2025} L.~Napolitano, Former Celsius Network CEO Alex Mashinsky sentenced to 12 years in prison, \emph{Decrypt}, 8 May 2025.

\bibitem{devries2021} A.~de~Vries and C.~Stoll, Bitcoin's growing e-waste problem, \emph{Resources, Conservation and Recycling}, 2021.

\bibitem{ethereum2022} Ethereum Foundation, \emph{The Merge}, ethereum.org, 2022.

\bibitem{flyxion2026ledger} Flyxion, \emph{Ledger Without Value: Contribution, Identity, and Governance after Tokenization}, essay, August 2026.

\bibitem{grossman1980} S.~J. Grossman and J.~E. Stiglitz, On the impossibility of informationally efficient markets, \emph{American Economic Review} 70(3), 393--408, 1980.

\bibitem{hirshleifer1971} J.~Hirshleifer, The private and social value of information and the reward to inventive activity, \emph{American Economic Review} 61(4), 561--574, 1971.

\bibitem{kyle1985} A.~S. Kyle, Continuous auctions and insider trading, \emph{Econometrica} 53(6), 1315--1335, 1985.

\bibitem{mckenzie2026} B.~McKenzie (director), \emph{Everyone Is Lying to You for Money}, CBC Docs, 2026.

\bibitem{nakamoto2008} S.~Nakamoto, \emph{Bitcoin: A peer-to-peer electronic cash system}, 2008.

\bibitem{nbc2025} Celsius CEO Alex Mashinsky sentenced to 12 years in multibillion-dollar crypto fraud case, \emph{NBC News}, 8 May 2025.

\bibitem{posner1975} R.~A. Posner, The social costs of monopoly and regulation, \emph{Journal of Political Economy} 83(4), 807--827, 1975.

\bibitem{tullock1967} G.~Tullock, The welfare costs of tariffs, monopolies, and theft, \emph{Western Economic Journal} 5(3), 224--232, 1967.

\end{thebibliography}

\end{document}
